
% ============================================================================
% DOCUMENT 5: California Code of Regulations, Title 22 - Physical AI Oncology Trial Site Regulations
% ============================================================================

\part{California Code of Regulations, Title 22 - Physical AI Oncology Trial Site Regulations}
\setcounter{section}{0}

\section{General Provisions}

\subsection{\S~75000. Authority and Purpose}

\begin{enumerate}[label=(\alph*)]
\item These regulations are adopted pursuant to the authority granted to the
  California Department of Public Health (hereinafter ``the department'') under
  the California Physical AI Oncology Clinical Trial Authorization and Site
  Establishment Act of 2026 (Health and Safety Code \S~1399.900 et seq.).

\item The purpose of these regulations is to establish standards and procedures
  for the authorization, inspection, and oversight of Physical AI oncology
  clinical trial sites throughout California.

\item These regulations apply uniformly across the State of California. The
  first authorized site shall be established in San Francisco, and subsequent
  sites may be authorized in any California jurisdiction that meets the
  requirements of these regulations.
\end{enumerate}

\subsection{\S~75001. Definitions}

In addition to the definitions established in Health and Safety Code
\S~1399.901, the following definitions apply:

\begin{enumerate}[label=(\alph*)]
\item ``Authorization certificate'' means the document issued by the department
  to a qualified entity permitting the operation of a Physical AI oncology
  clinical trial site.

\item ``Graduated autonomy'' means the phased approach to Physical AI system
  deployment in which autonomy levels increase from teleoperation in Phase 1
  to supervised autonomy in Phase 2 to full autonomy range in Phase 3, as
  established in the adapted 21 CFR Part 312 \S~312.21~\cite{cfr312adapt}.

\item ``Pre-procedure safety matrix'' means the checklist of safety
  verifications that shall be completed before initiating any Physical
  AI-assisted clinical procedure, as defined in the adapted 21 CFR Part 50
  \S~50.3~\cite{cfr50adapt}.

\item ``Robot capability profile'' means the formal documentation of a Physical
  AI system's validated task capabilities, performance envelope, and safety
  constraints.

\item ``Site safety officer'' means a human staff member certified by the
  department to supervise Physical AI system operations during a shift, with
  authority to initiate emergency stop procedures and override autonomous
  operations.
\end{enumerate}

\section{Authorization Process}

\subsection{\S~75010. Application Requirements}

\begin{enumerate}[label=(\alph*)]
\item An entity seeking authorization to operate a Physical AI oncology
  clinical trial site shall submit an application to the department on forms
  prescribed by the department, including the following:
  \begin{enumerate}[label=(\arabic*)]
  \item Identification of the applicant entity, including ownership structure,
    officers, and any parent organizations.
  \item Site location and facility description, including architectural plans,
    engineering specifications, and evidence of compliance with applicable
    building codes.
  \item Physical AI System Deployment Plan identifying all robot types and
    instances, with current USL scores for each system~\cite{usl2026}.
  \item Phase 0 simulation validation results for each Physical AI system
    type, demonstrating cross-framework validation using two or more
    frameworks~\cite{iche6r3adapt}.
  \item Staffing plan identifying all positions, qualifications, training
    requirements, and shift schedules for the 24-hour operating model.
  \item Quality management system documentation addressing Physical AI-specific
    processes.
  \item Emergency response plan addressing Physical AI system failures,
    cybersecurity incidents, natural disasters, and medical emergencies.
  \item Patient flow projections based on the simulation data showing capacity
    for 150 to 180 patients per day across 15 or more cancer
    types~\cite{newtrial2026}.
  \item Financial documentation demonstrating adequate resources to sustain
    operations for a minimum of 24 months.
  \item Evidence of compliance with adapted ICH E6(R3), 21 CFR Part 50, and
    21 CFR Part 312 regulatory frameworks~\cite{iche6r3adapt,cfr50adapt,cfr312adapt}.
  \end{enumerate}

\item The application shall be accompanied by a nonrefundable filing fee of
  twenty-five thousand dollars (\$25,000).
\end{enumerate}

\subsection{\S~75011. Review Process}

\begin{enumerate}[label=(\alph*)]
\item The department shall assign a review team consisting of specialists in
  robotic engineering, oncology, clinical trial regulation, facility safety,
  and cybersecurity.

\item The review team shall complete an initial completeness review within 30
  days of application receipt and notify the applicant of any deficiencies.

\item The department shall conduct a substantive review of the complete
  application and issue a decision within 120 days of receiving a complete
  application.

\item The department may conduct a pre-authorization site inspection as part
  of the review process.
\end{enumerate}

\subsection{\S~75012. Authorization Certificate}

\begin{enumerate}[label=(\alph*)]
\item An authorization certificate shall specify the following:
  \begin{enumerate}[label=(\arabic*)]
  \item The authorized Physical AI system types and maximum number of instances.
  \item The maximum number of concurrent patients.
  \item The authorized operating schedule.
  \item The minimum staffing requirements.
  \item Any conditions or limitations specific to the site.
  \end{enumerate}

\item The authorization certificate shall be valid for a period of three years,
  subject to annual compliance verification.

\item The certificate shall be prominently displayed at the site entrance and
  available for public inspection.
\end{enumerate}

\section{Physical AI System Standards}

\subsection{\S~75020. System Registration}

\begin{enumerate}[label=(\alph*)]
\item Each Physical AI system instance deployed at an authorized site shall be
  registered with the department, including the following information:
  \begin{enumerate}[label=(\arabic*)]
  \item System type, manufacturer, model, and serial number.
  \item Current USL score with dimension-level breakdown (simulation framework
    switching, AI integration, cross-robot progress sharing, clinical trial
    collaboration)~\cite{usl2026}.
  \item Current PSL scores across three regulatory dimensions (Omniscient,
    Omnipresent, Omnipotent)~\cite{pslframework2026}.
  \item Robot capability profile documenting validated task capabilities.
  \item Phase 0 simulation validation results.
  \item Installation date and commissioning records.
  \end{enumerate}

\item Registration shall be updated within 14 days of any change to system
  configuration, USL score, or capability profile.
\end{enumerate}

\subsection{\S~75021. Minimum USL Thresholds}

Physical AI systems shall meet the following minimum USL score thresholds,
consistent with the MCP-PAI standard~\cite{cfr50adapt}:

\begin{center}
\begin{tabular}{lc}
\toprule
\textbf{System Category} & \textbf{Minimum USL Score} \\
\midrule
Surgical robots & 7.0 \\
Radiotherapy positioning robots & 6.0 \\
Robotic needle-placement systems & 6.0 \\
Radiotherapy motion-tracking robots & 6.0 \\
Imaging assistant robots & 6.0 \\
Steerable needle robots & 6.0 \\
Collaborative robots (cobots) & 5.0 \\
Rehabilitation exoskeletons & 4.0 \\
Humanoid robots & 4.0 \\
Social companion robots & 3.0 \\
\bottomrule
\end{tabular}
\end{center}

\subsection{\S~75022. Maintenance and Calibration}

\begin{enumerate}[label=(\alph*)]
\item Each Physical AI system shall undergo scheduled maintenance at intervals
  specified by the manufacturer and documented in the robot capability profile.

\item Maintenance windows shall be coordinated to maintain continuous clinical
  capability across all system categories, as demonstrated in the simulation
  where SURG-01 maintenance (22:15 to 01:59) was scheduled during low-demand
  hours with SURG-02 and SURG-03 remaining
  available~\cite{newtrial2026}.

\item Calibration records shall include date, calibration parameters, results,
  and the identity of the person or system performing the calibration.

\item A Physical AI system shall not be used for patient procedures until
  post-maintenance verification demonstrates performance within specified
  tolerances.
\end{enumerate}

\section{Staffing and Training Standards}

\subsection{\S~75030. Minimum Staffing Requirements}

\begin{enumerate}[label=(\alph*)]
\item An authorized Physical AI oncology clinical trial site operating on a
  24-hour schedule shall maintain the following minimum staffing:
  \begin{enumerate}[label=(\arabic*)]
  \item Two to three site safety officers per shift, with three shifts covering
    the 24-hour operating period.
  \item One on-call emergency physician available within 15 minutes at all
    times.
  \item One on-call pharmacist available within 15 minutes at all times.
  \item One facility maintenance technician on-site per shift.
  \item One cybersecurity operations specialist on-call at all times.
  \end{enumerate}

\item Total full-time equivalent staffing shall be a minimum of 8 and a maximum
  of 10, reflecting the Physical AI model in which more patients are treated
  with more robots and fewer workers~\cite{sitespec2026}.

\item The site shall maintain a staffing contingency plan for situations
  requiring additional human personnel, such as multiple simultaneous
  adverse events or extended system outages.
\end{enumerate}

\subsection{\S~75031. Training and Certification}

\begin{enumerate}[label=(\alph*)]
\item Site safety officers shall complete a department-approved training program
  of not less than 160 hours covering the following:
  \begin{enumerate}[label=(\arabic*)]
  \item Physical AI system operation and capabilities for all 10 robot
    categories deployed at the site.
  \item Emergency stop procedures and manual override protocols.
  \item Adverse event recognition, documentation, and reporting.
  \item Patient interaction protocols, including patient rights under AB 2847.
  \item Cybersecurity awareness and incident response.
  \item Regulatory compliance including ICH E6(R3), 21 CFR Part 50, and
    21 CFR Part 312 as adapted for Physical AI~\cite{iche6r3adapt,cfr50adapt,cfr312adapt}.
  \end{enumerate}

\item Certification shall be valid for two years, with 40 hours of continuing
  education required for renewal.

\item The department shall develop and maintain the training curriculum in
  consultation with Physical AI system manufacturers, oncology professional
  organizations, and patient advocacy groups.
\end{enumerate}

\section{Inspection and Enforcement}

\subsection{\S~75040. Routine Inspections}

\begin{enumerate}[label=(\alph*)]
\item The department shall conduct routine inspections of authorized sites at
  least quarterly, without prior notice.

\item Inspections shall assess compliance with authorization conditions,
  Physical AI system registration and maintenance, staffing requirements,
  patient rights protections, adverse event reporting, data protection, and
  facility condition.

\item The department shall publish inspection reports on its website within 30
  days of the inspection.
\end{enumerate}

\subsection{\S~75041. Enforcement Actions}

\begin{enumerate}[label=(\alph*)]
\item The department may take the following enforcement actions for
  noncompliance:
  \begin{enumerate}[label=(\arabic*)]
  \item Written notice of deficiency with required corrective action plan.
  \item Administrative penalty of up to ten thousand dollars (\$10,000) per
    day per violation.
  \item Suspension of authorization to use specific Physical AI system types.
  \item Suspension of site authorization.
  \item Revocation of site authorization.
  \end{enumerate}

\item Before suspending or revoking authorization, the department shall provide
  written notice and an opportunity for hearing in accordance with the
  Administrative Procedure Act (Government Code \S~11500 et seq.).

\item In cases of imminent danger to patient safety, the department may issue
  an immediate suspension order, with a hearing to be held within 15 days.
\end{enumerate}

\section{Reporting Requirements}

\subsection{\S~75050. Monthly Reports}

\begin{enumerate}[label=(\alph*)]
\item Authorized sites shall submit monthly reports to the department including:
  \begin{enumerate}[label=(\arabic*)]
  \item Total patients served and procedures performed, disaggregated by
    cancer type and Physical AI system type.
  \item Current PSL scores for all robot types and cumulative site PSL score.
  \item All adverse events with grade, resolution, and root cause analysis.
  \item Physical AI system utilization rates and maintenance records.
  \item Staffing levels and any departures from minimum requirements.
  \end{enumerate}
\end{enumerate}

\subsection{\S~75051. Annual Reports}

\begin{enumerate}[label=(\alph*)]
\item Authorized sites shall submit annual reports including all elements
  required under the adapted 21 CFR Part 312 \S~312.33, plus the following
  state-specific elements~\cite{cfr312adapt}:
  \begin{enumerate}[label=(\arabic*)]
  \item Year-over-year comparison of patient outcomes and throughput.
  \item Updated USL scores for all Physical AI systems with trend analysis.
  \item Community benefits achieved, including local workforce development.
  \item Patient satisfaction data and complaint resolution statistics.
  \item Comparison of site performance to the baseline established by the
    24-hour simulation (168 patients, 99.7 percent uptime, 8-minute average
    wait, 4.7 out of 5.0 satisfaction)~\cite{newtrial2026}.
  \end{enumerate}
\end{enumerate}

